\documentclass[11pt,a4paper]{article}

% --- PACOTES FUNDAMENTAIS ---
\usepackage[utf8]{inputenc}
\usepackage[T1]{fontenc}
\usepackage[brazil]{babel}
\usepackage{geometry}
\usepackage{graphicx}
\usepackage{fancyhdr}
\usepackage{lastpage}
\usepackage{titlesec}
\usepackage{booktabs}
\usepackage{tabularx}
\usepackage{xcolor}
\usepackage{float}
\usepackage{caption}
\usepackage{subcaption}
\usepackage{hyperref}

% --- CONFIGURAÇÕES DE PÁGINA E MARGENS ---
\geometry{
    top=2.5cm,
    bottom=2.5cm,
    left=2.5cm,
    right=2.5cm
}

% --- CORES FORENSES / IDENTIDADE VISUAL ---
\definecolor{primary}{RGB}{31, 32, 34}       % Grafite institucional
\definecolor{accent}{RGB}{184, 134, 11}      % Dourado clássico (Dark Goldenrod)
\definecolor{lightgray}{RGB}{245, 245, 245}  % Fundo suave de tabelas
\definecolor{bordergray}{RGB}{210, 210, 210} % Linhas de tabela

% --- CABEÇALHO E RODAPÉ ---
\pagestyle{fancy}
\fancyhf{}
\fancyhead[L]{\small\textbf{\color{primary}LAUDO DE VISTORIA PERICIAL} — \textit{Consumo de Energia}}
\fancyhead[R]{\small\color{gray}Proc. N.º {{NUMERO_PROCESSO}}}
\fancyfoot[L]{\footnotesize\color{gray}VistoriaPro — Sistema de Engenharia Diagnóstica}
\fancyfoot[R]{\footnotesize\color{gray}Página \thepage\ de \pageref{LastPage}}
\renewcommand{\headrulewidth}{0.5pt}
\renewcommand{\footrulewidth}{0.5pt}

% --- FORMATAÇÃO DE SEÇÕES ---
\titleformat{\section}
  {\normalfont\large\bfseries\color{primary}}
  {\thesection.}{0.5em}{}
  [{\color{accent}\titlerule[1.5pt]}]

\titleformat{\subsection}
  {\normalfont\normalsize\bfseries\color{primary}}
  {\thesubsection}{0.5em}{}

\setlength{\parindent}{0pt}
\setlength{\parskip}{6pt}

\begin{document}

% --- CAPA / CABEÇALHO DO DOCUMENTO ---
\begin{center}
    {\Large\textbf{\color{primary}LAUDO TÉCNICO DE CONSTATAÇÃO PERICIAL}}\\[4pt]
    {\normalsize\textbf{\color{accent}PERÍCIA JUDICIAL EM ENERGIA ELÉTRICA}}\\[12pt]
    \rule{\textwidth}{0.8pt}
\end{center}

\vspace{0.2cm}

% --- QUADRO DE IDENTIFICAÇÃO DO PROCESSO ---
\begin{table}[H]
\centering
\begin{tabularx}{\textwidth}{@{}l X@{}}
\toprule
\textbf{Processo Judicial N.º:} & {{NUMERO_PROCESSO}} \\
\textbf{Tipo de Ação:}        & {{TIPO_ACAO}} \\
\textbf{Autor(a):}            & \textbf{{{NOME_AUTOR}}} \\
\textbf{Réu / Concessionária:} & {{REU_CONCESSIONARIA}} \\
\textbf{Data da Diligência:}  & {{DATA_VISTORIA}} \\
\textbf{N.º da Vistoria / Período:} & Vistoria {{NUMERO_VISTORIA}} — {{PERIODO_VISTORIA}} \\
\bottomrule
\end{tabularx}
\end{table}

\vspace{0.4cm}

% --- SEÇÃO 1: OBJETIVO E LOCAL DA DILIGÊNCIA ---
\section{Identificação da Diligência e Presença das Partes}

Na data e horário supramencionados, procedeu-se à vistoria técnica presencial no imóvel objeto da presente lide pericial. Durante os trabalhos de campo, foram constatadas as seguintes presenças:

\begin{itemize}
    \item \textbf{Representação do Autor:} {{REPRESENTACAO_AUTOR}}
    \item \textbf{Representação do Réu / Concessionária:} {{REPRESENTACAO_REU}}
\end{itemize}

\textit{{{OBSERVACOES_PRESENCA}}}

\vspace{0.4cm}

% --- SEÇÃO 2: DETALHES DO MEDIDOR E CARACTERÍSTICAS TÉCNICAS ---
\section{Constatações da Unidade Consumidora e Medição}

Procedeu-se à inspeção minuciosa dos equipamentos de medição instalados na unidade consumidora:

\begin{table}[H]
\centering
\begin{tabularx}{\textwidth}{@{}X X@{}}
\toprule
\textbf{Parâmetro Inspecionado} & \textbf{Constatação Pericial} \\
\midrule
Número do Medidor               & \textbf{{{NUMERO_MEDIDOR}}} \\
Tecnologia / Chip Integrado     & {{MEDIDOR_CHIP}} \\
Condições Físicas e Lacres      & {{CONDICOES_MEDIDOR}} \\
Fornecimento Ativo no Momento   & {{CORTE_ENERGIA}} \\
Total de Moradores / Residentes & {{QTD_PESSOAS}} pessoa(s) \\
Quantidade de Cômodos do Imóvel & {{QTD_COMODOS}} cômodo(s) \\
\bottomrule
\end{tabularx}
\end{table}

\vspace{0.4cm}

% --- SEÇÃO 3: LEVANTAMENTO DE CARGAS E ELETRODOMÉSTICOS ---
\section{Levantamento de Cargas e Equipamentos Elétricos}

Realizou-se o inventário de todos os equipamentos elétricos existentes no imóvel no momento da perícia para composição da curva de demanda estimada:

{{TABELA_ELETRODOMESTICOS}}

\vspace{0.4cm}

% --- SEÇÃO 4: CHECKLIST TÉCNICO ---
\section{Checklist de Conformidade Técnica}

\begin{itemize}
{{CHECKLIST_ITENS}}
\end{itemize}

\vspace{0.4cm}

% --- SEÇÃO 5: OBSERVAÇÕES E CONCLUSÕES TÉCNICAS ---
\section{Observações Finais do Perito}

{{OBSERVACOES_FINAIS}}

\vspace{0.8cm}

% --- SEÇÃO 6: REGISTRO FOTOGRÁFICO (ANEXO DE EVIDÊNCIAS) ---
\newpage
\section{Anexo Fotográfico de Evidências}
\label{sec:fotos}

As fotografias a seguir foram capturadas in loco pelo perito durante a diligência técnica e representam fielmente o estado das instalações:

\subsection{Fotografias do Imóvel e Fachada}
{{FOTOS_IMOVEL_LATEX}}

\vspace{0.6cm}

\subsection{Fotografias do Medidor e Detalhes Técnicos}
{{FOTOS_MEDIDOR_LATEX}}

\vspace{1.5cm}

% --- ASSINATURA DO PERITO ---
\begin{center}
    \rule{8cm}{0.6pt}\\[4pt]
    \textbf{Fabíola — Arquiteta \& Engenharia Diagnóstica}\\[2pt]
    Perita Judicial Especialista\\
    \textit{Documento gerado e autenticado via VistoriaPro}
\end{center}

\end{document}
